\documentclass[a4paper,12pt]{article}
\usepackage[T2A]{fontenc}
\usepackage[utf8]{inputenc}
\usepackage[russian]{babel}
\usepackage{amsmath,amssymb,mathtools}
\usepackage{PTSans}
\renewcommand{\familydefault}{\sfdefault}
\usepackage[margin=23mm,top=23mm,bottom=24mm,headheight=25pt]{geometry}
\usepackage{microtype}
\usepackage{tikz}
\usetikzlibrary{arrows.meta,positioning,shapes.geometric,calc}
\usepackage{tabularx,booktabs}
\usepackage{enumitem}
\usepackage{parskip}
\usepackage{fancyhdr}
\usepackage{setspace}
\usepackage{xcolor}
\usepackage{titlesec}
\usepackage{hyperref}
\usepackage{needspace}
\definecolor{accent}{HTML}{087A7C}
\definecolor{darkaccent}{HTML}{15575B}
\definecolor{soft}{HTML}{667276}
\definecolor{line}{HTML}{D5E2E0}
\definecolor{ink}{HTML}{252D2D}
\color{ink}
\hypersetup{colorlinks=true,linkcolor=darkaccent,urlcolor=darkaccent}
\setstretch{1.13}
\setlength{\parindent}{0pt}
\setlength{\parskip}{5pt}
\setlist[itemize]{leftmargin=6mm,topsep=3pt,itemsep=3pt}
\setlist[enumerate]{leftmargin=7mm,topsep=3pt,itemsep=5pt}
\titleformat{\section}{\Large\bfseries\color{darkaccent}}{\thesection}{0.7em}{}
\titleformat{\subsection}{\large\bfseries\color{darkaccent}}{\thesubsection}{0.6em}{}
\titlespacing*{\section}{0pt}{1.1em}{0.55em}
\titlespacing*{\subsection}{0pt}{0.85em}{0.28em}
\pagestyle{fancy}
\fancyhf{}
\fancyhead[L]{\small\color{soft} Математика \quad\textcolor{line}{|}\quad Движение по реке}
\fancyhead[R]{\small\color{soft} Ярослав Верещагин}
\fancyfoot[L]{\footnotesize\color{soft} 09.11.2026}
\fancyfoot[R]{\footnotesize\color{soft} \thepage}
\renewcommand{\headrulewidth}{0.35pt}
\renewcommand{\headrule}{\hbox to\headwidth{\color{line}\leaders\hrule height \headrulewidth\hfill}}
\newcommand{\important}[1]{\textcolor{darkaccent}{\textbf{#1}}}
\newcommand{\task}[1]{\medskip\noindent\important{#1}\par}
\newcommand{\checkline}[1]{\item[$\square$] #1}
\begin{document}
\begin{titlepage}
\thispagestyle{empty}
\begin{center}
\vspace*{23mm}
{\small\bfseries\color{accent}\MakeUppercase{Математика}}\par
\vspace{24mm}
{\fontsize{28}{33}\selectfont\bfseries\color{darkaccent} Чек-лист}\par
\vspace{9mm}
{\fontsize{22}{27}\selectfont\bfseries Движение по реке}\par
\vspace{4mm}
{\large Скорость, время, расстояние. Рейсы, стоянки и задержки}\par
\vspace{23mm}
{\normalsize Подготовка к ОГЭ по математике}\par
\vspace{20mm}
{\Large\bfseries Ярослав Верещагин}\par
\vspace{8mm}
{\normalsize 9 ноября 2026 г.}\par
\vfill
{\normalsize\color{soft} Лёвин Артём Александрович}\par
{\normalsize\color{soft} эксклюзивно для Ярослава Верещагина}\par
\vspace{22mm}
\end{center}
\begin{tikzpicture}[remember picture,overlay]
\draw[accent!45,line width=1.2pt] ([xshift=9mm,yshift=15mm]current page.south west)
.. controls ([xshift=62mm,yshift=35mm]current page.south west) and
([xshift=-67mm,yshift=4mm]current page.south east) ..
([xshift=-9mm,yshift=22mm]current page.south east);
\end{tikzpicture}
\end{titlepage}
\setcounter{page}{1}

\section{Главная идея: три величины}
\begin{itemize}[label={$\square$}]
\checkline{Я различаю \important{расстояние} $S$ (км), \important{скорость} $v$ (км/ч) и \important{время} $t$ (ч).}
\checkline{Помню основную формулу и умею выразить любую величину:}
\end{itemize}
\[
\boxed{S=vt}, \qquad \boxed{t=\frac{S}{v}}, \qquad \boxed{v=\frac{S}{t}}.
\]
\noindent В задаче о реке всегда нужно определить, какую скорость относительно берега имеет судно. Течение помогает двигаться вниз по реке и замедляет движение вверх.

\subsection{Скорость относительно берега}
Пусть $x$~--- скорость лодки в неподвижной воде, а $u$~--- скорость течения.
\[
\boxed{v_{\text{по теч.}}=x+u}, \qquad
\boxed{v_{\text{против теч.}}=x-u},\qquad x>u>0.
\]
\begin{center}
\begin{tikzpicture}[>=Latex,scale=1]
\draw[thick,gray!60] (0,0)--(11.6,0);
\fill[accent!6] (0,-0.30) rectangle (11.6,0.30);
\node[above] at (0.45,0.28) {$A$};
\node[above] at (11.15,0.28) {$B$};
\draw[accent,very thick,->] (1.25,0.16)--(10.25,0.16)
 node[midway,above=1pt,fill=white,inner sep=1pt]{\footnotesize по течению: $x+u$};
\draw[orange!80!black,very thick,->] (10.25,-0.16)--(1.25,-0.16)
 node[midway,below=1pt,fill=white,inner sep=1pt]{\footnotesize против течения: $x-u$};
\end{tikzpicture}
\end{center}
\begin{itemize}[label={$\square$}]
\checkline{Понимаю, почему скорость \important{плота}, который свободно плывёт с течением, равна $u$.}
\checkline{Понимаю, что на участках $A\to B$ и $B\to A$ одинаковое расстояние, если лодка возвращается той же дорогой.}
\checkline{Помню: для деления на $x-u$ необходимо условие $x>u$.}
\end{itemize}

\subsection{Как оформить таблицу}
\renewcommand{\arraystretch}{1.22}
\begin{center}
\begin{tabularx}{\linewidth}{@{}Xccc@{}}
\toprule
Участок пути & Скорость & Время & Расстояние \\
\midrule
Вниз по реке & $x+u$ & $\dfrac{S}{x+u}$ & $S$ \\
Вверх по реке & $x-u$ & $\dfrac{S}{x-u}$ & $S$ \\
Стоянка у пристани & $0$ & $t_{\text{ст}}$ & $0$ \\
Плот по течению & $u$ & $\dfrac{S_{\text{пл}}}{u}$ & $S_{\text{пл}}$ \\
\bottomrule
\end{tabularx}
\end{center}

\section{Как переводить слова в уравнение}
\important{Разность времени.} Если возвращение по течению заняло на $3$ часа меньше, пишем $t_{\text{против}}-t_{\text{по}}=3$.

\important{Время со стоянкой.} Общая длительность равна сумме времён двух переходов и стоянки: $t_{\text{общ}}=t_{\text{туда}}+t_{\text{ст}}+t_{\text{обратно}}$.

\important{Задержка отправления.} Если плот вышел на $3$ часа раньше лодки, до общего момента наблюдения он плыл на $3$ часа дольше. Для расписания сначала находим длительность: от 10:00 до 21:00 проходит $11$ часов.

\section{Четыре разобранных случая}
\subsection{1. Время движения вверх и вниз различается}
\textbf{Задача Т4.2.} Баржа проходит $24$ км против течения и возвращается обратно по тому же пути. Обратная дорога занимает на $3$ часа меньше. Скорость течения $2$ км/ч. Найдём скорость баржи в неподвижной воде.

Обозначим её через $x$ км/ч. Тогда $x>2$. Вверх по реке скорость равна $x-2$, вниз~--- $x+2$. Расстояния одинаковые, поэтому:
\[
\frac{24}{x-2}-\frac{24}{x+2}=3.
\]
Разделим обе части на $3$, приведём к общему знаменателю и сократим числитель:
\[
\frac{8}{x-2}-\frac{8}{x+2}=1,\qquad
\frac{8(x+2)-8(x-2)}{(x-2)(x+2)}=1,
\]
\[
\frac{32}{x^2-4}=1,\qquad x^2=36.
\]
Корни $x=6$ и $x=-6$. Условию $x>2$ соответствует только $x=6$.
\textbf{Проверка:} против течения $24:4=6$ ч, по течению $24:8=3$ ч. Разница действительно $3$ ч. \important{Ответ: $6$ км/ч.}

\subsection{2. В суммарное время входит стоянка}
\textbf{Задача Т4.3.} Теплоход со скоростью $18$ км/ч в неподвижной воде плывёт вниз по течению, стоит $3$ часа и возвращается назад. Течение имеет скорость $2$ км/ч. Весь рейс занимает $12$ часов. Сколько километров прошёл теплоход за оба перехода?

Пусть $s$ км~--- расстояние в \emph{одну} сторону. Тогда скорости равны $20$ и $16$ км/ч. Записываем все три промежутка времени:
\[
\frac{s}{20}+3+\frac{s}{16}=12,
\qquad \frac{s}{20}+\frac{s}{16}=9.
\]
Умножая на $80$, получаем $4s+5s=720$, откуда $s=80$ км. Весь путь состоит из двух одинаковых по длине участков: $80+80=160$ км.
\textbf{Проверка:} $4+3+5=12$ ч. \important{Ответ: $160$ км.}

\Needspace{16\baselineskip}
\subsection{3. Время задано по часам}
\textbf{Задача Т4.4.} От пристани $A$ до $B$~--- $72$ км. Теплоход уходит из $A$ в 10:00, стоит в $B$ четыре часа и возвращается в $A$ в 21:00. Скорость течения $3$ км/ч. Найдём скорость теплохода в неподвижной воде.

От 10:00 до 21:00 проходит $11$ часов. После исключения стоянки на оба перехода приходится $11-4=7$ часов. Пусть искомая скорость~--- $x>3$ км/ч. Тогда
\[
\frac{72}{x+3}+\frac{72}{x-3}=7.
\]
После умножения на $(x+3)(x-3)$ получаем
\[
72(x-3)+72(x+3)=7(x^2-9),
\]
\[
7x^2-144x-63=0.
\]
Дискриминант: $D=144^2+4\cdot7\cdot63=22500=150^2$. Поэтому
\[
x=\frac{144\pm150}{14}.
\]
Корень $x=21$ удовлетворяет $x>3$; второй корень отрицательный.
\textbf{Проверка:} $72:24=3$ ч, $72:18=4$ ч, стоянка $4$ ч; всего $11$ ч. \important{Ответ: $21$ км/ч.}

\subsection{4. Плот стартовал раньше моторной лодки}
\textbf{Задача Т4.5.} Между пристанями $A$ и $B$~--- $72$ км. Плот отплывает из $A$ вниз по реке. Спустя $3$ часа за ним отправляется моторная лодка; она доплывает до $B$ и сразу возвращается в $A$. К её возвращению плот успел пройти $39$ км. Скорость течения $3$ км/ч.

Плот плывёт $39:3=13$ часов. Лодка начала движение на $3$ часа позже, поэтому её полный рейс занял $13-3=10$ часов. При скорости лодки $x>3$ км/ч:
\[
\frac{72}{x+3}+\frac{72}{x-3}=10.
\]
Умножаем на $x^2-9$ и решаем квадратное уравнение:
\[
144x=10(x^2-9),\quad 5x^2-72x-45=0,
\]
\[
D=72^2+4\cdot5\cdot45=6084=78^2,
\qquad x=\frac{72\pm78}{10}.
\]
Допустимый корень $x=15$. \textbf{Проверка:} лодка потратила $72:18=4$ ч и $72:12=6$ ч; плот двигался ещё $3$ ч до её старта, итого $13$ ч, или $39$ км.
\important{Ответ: $15$ км/ч.}

\clearpage
\section*{Алгоритм решения задачи о движении по реке}
\vspace{8mm}
\begin{center}
\begin{tikzpicture}[
>=Latex,
node distance=11mm,
startstop/.style={ellipse,draw=darkaccent,fill=accent!5,minimum width=66mm,minimum height=11mm,align=center,font=\small},
action/.style={rectangle,rounded corners=2pt,draw=accent!70,fill=accent!3,text width=76mm,minimum height=13mm,align=center,font=\small,inner sep=4pt},
choice/.style={diamond,aspect=2.5,draw=darkaccent,fill=accent!3,inner sep=3pt,align=center,font=\small,text width=43mm},
arrow/.style={->,thick,color=darkaccent}
]
\node[startstop] (a) {Начало: что известно и что ищем?};
\node[action,below=of a] (b) {Обозначьте неизвестную: $x$ или $s$. Запишите условия: $x>u>0$, $s>0$.};
\node[action,below=of b] (c) {Для каждого участка найдите скорость относительно берега, затем время $t=S/v$.};
\node[choice,below=13mm of c] (d) {Есть стоянка или задержка старта?};
\node[action,right=15mm of d,text width=41mm] (e) {Учтите дополнительное время в общей длительности.};
\node[action,below=18mm of d] (f) {Составьте уравнение по сумме или разности промежутков времени.};
\node[action,below=of f] (g) {Решите уравнение, проверьте допустимость корней.};
\node[startstop,below=of g] (h) {Проверьте условие и запишите ответ с единицами.};
\draw[arrow] (a)--(b);
\draw[arrow] (b)--(c);
\draw[arrow] (c)--(d);
\draw[arrow] (d)--node[left,font=\footnotesize]{нет}(f);
\draw[arrow] (d)--node[above,font=\footnotesize]{да}(e);
\draw[arrow] (e.south) -- ++(0,-8mm) -| (f.east);
\draw[arrow] (f)--(g);
\draw[arrow] (g)--(h);
\end{tikzpicture}
\end{center}

\clearpage
\section{Типичные ошибки и самопроверка}
\begin{itemize}[label={$\square$}]
\checkline{При движении против течения использую $x-u$ и отдельно проверяю $x>u$.}
\checkline{Если требуется \emph{весь путь}, складываю длины всех переходов. В задачах с возвращением это часто $2S$.}
\checkline{Если есть стоянка, добавляю её ко времени рейса; пройденное за стоянку расстояние равно нулю.}
\checkline{Если плот вышел раньше, продолжительность движения плота больше продолжительности движения лодки на величину задержки.}
\checkline{Если условие сообщает время прибытия и отправления, сначала нахожу длительность в часах.}
\checkline{После решения квадратного уравнения исключаю корни, несовместимые с физическим смыслом.}
\checkline{Проверяю результат в исходном условии, подставляя найденную скорость в расчёт времени.}
\end{itemize}
\vspace{3mm}
\noindent\important{Контрольный вопрос для себя:} понимаю ли я, почему в уравнение попали именно эти слагаемые и что обозначает каждая дробь? При затруднении восстановите таблицу «скорость~--- время~--- расстояние».

\clearpage
\section*{Тренировочный блок}
\textcolor{soft}{Путь А. 10 заданий, расположенных по возрастанию сложности. После каждого задания запишите таблицу движения или краткую схему времени. Решения приведите самостоятельно; ответы находятся в конце пособия.}

\subsection*{Базовый уровень}
\begin{enumerate}[label=\textcolor{accent}{\bfseries\arabic*.},leftmargin=8mm,itemsep=8pt]
\item \textbf{Прогулочный катер.} От пристани $A$ до пристани $B$ по течению реки $48$ км. Скорость прогулочного катера в неподвижной воде равна $14$ км/ч, а скорость течения~--- $2$ км/ч. Катер без остановок добрался до $B$, сразу развернулся и вернулся в $A$ тем же маршрутом. Найдите скорость катера на каждом участке, время движения туда и обратно, а также общую продолжительность прогулки.

\item \textbf{Путешествие плота.} Утром в 8:30 туристы отпустили небольшой плот, который свободно двигался вместе с течением реки. Скорость течения на этом участке постоянна и равна $3$ км/ч. Наблюдатели встретили плот ниже по реке, когда он преодолел $27$ км. Сколько часов плот находился в пути? Во сколько он оказался в месте встречи? Объясните, почему для плота не требуется вводить отдельную скорость в неподвижной воде.

\item \textbf{Две измеренные скорости.} На одном и том же участке реки лодка с работающим мотором шла вниз по течению со скоростью $18$ км/ч, а при возвращении вверх по течению~--- со скоростью $12$ км/ч. Мощность мотора оставалась одинаковой, и скорость течения не изменялась. Найдите скорость лодки в неподвижной воде и скорость течения. Проверьте, что полученные значения дают обе наблюдаемые скорости.
\end{enumerate}

\subsection*{Средний уровень}
\begin{enumerate}[label=\textcolor{accent}{\bfseries\arabic*.},leftmargin=8mm,itemsep=8pt,start=4]
\item \textbf{Возвращение быстрее.} Моторная лодка прошла от пристани $A$ до пристани $B$ $30$ км против течения. После короткого разворота она вернулась в $A$ по течению без остановок. Обратный переход оказался на $2$ часа короче первого. Скорость речного течения составляет $2$ км/ч. Определите скорость лодки в неподвижной воде. Составляя уравнение, отдельно укажите, почему время движения вверх больше времени движения вниз.

\item \textbf{Определение скорости течения.} Катер способен двигаться со скоростью $10$ км/ч в неподвижной воде. Он прошёл $48$ км вниз по реке, затем столько же вверх. Движение против течения заняло на $2$ часа больше, чем движение по течению. Найдите скорость течения, считая её постоянной. Проверьте найденное значение по продолжительности обоих переходов.

\end{enumerate}
\clearpage
\section*{Тренировочный блок: продолжение}
\subsection*{Средний уровень}
\begin{enumerate}[label=\textcolor{accent}{\bfseries\arabic*.},leftmargin=8mm,itemsep=8pt,start=6]
\item \textbf{Рейс с ожиданием у пристани.} Теплоход отправился по течению к соседнему городу, простоял там $2$ часа для высадки пассажиров и отправился в обратный путь против течения. Вся поездка, считая с момента отплытия до возвращения, продолжалась $11$ часов. Собственная скорость теплохода равна $18$ км/ч, скорость течения~--- $2$ км/ч. Найдите расстояние между городами и общий путь теплохода за два перехода. Обратите внимание, что во время стоянки теплоход не преодолевает расстояние.
\end{enumerate}

\subsection*{Уровень выше среднего}
\begin{enumerate}[label=\textcolor{accent}{\bfseries\arabic*.},leftmargin=8mm,itemsep=12pt,start=7]
\item \textbf{Поездка по расписанию.} В 7:45 теплоход отправился от пристани $A$ по течению к пристани $B$, расположенной в $80$ км. В пункте $B$ теплоход находился ровно $1$ час, после чего вернулся к пристани $A$ против течения. В тот же день в 17:45 он завершил рейс. Скорость течения реки равна $2$ км/ч. С какой скоростью теплоход движется в неподвижной воде? Перед составлением уравнения определите полное время рейса и продолжительность движения без стоянки.

\item \textbf{Плот и догоняющая его моторная лодка.} Из пункта $A$ вниз по реке отправился плот. Спустя $2$ часа из того же пункта вышла моторная лодка, прошла $48$ км по течению до пристани $B$, сразу повернула обратно и вернулась в $A$. К моменту возвращения лодки плот удалился от $A$ на $18$ км. Скорость течения составляет $2$ км/ч. Найдите скорость лодки в неподвижной воде. Особое внимание уделите тому, сколько времени двигался каждый участник и почему эти промежутки различаются.

\item \textbf{Неизвестная скорость реки.} Пассажирский теплоход со скоростью $15$ км/ч в неподвижной воде отправился по течению от пристани $A$ к пристани $B$, находящейся в $72$ км. В пункте $B$ он простоял $1$ час, затем вернулся обратно. От выхода из $A$ до возвращения туда прошло $11$ часов. Вычислите скорость течения реки, считая её постоянной. Проверьте, что теплоход способен идти против течения с найденной скоростью.
\end{enumerate}

\subsection*{Сложный уровень}
\begin{enumerate}[label=\textcolor{accent}{\bfseries\arabic*.},leftmargin=8mm,itemsep=10pt,start=10]
\item \textbf{Плот, задержка отправления и остановка теплохода.} С пристани $A$ по течению реки отправился свободно плывущий плот. Ровно через $2$ часа после него от той же пристани вышел моторный теплоход. Теплоход проплыл $80$ км вниз по реке до пристани $B$, простоял там $1$ час, затем преодолел те же $80$ км против течения и вернулся в $A$. К моменту его возвращения плот находился на расстоянии $24$ км от $A$. Скорость течения равна $2$ км/ч. Найдите скорость теплохода в неподвижной воде. Сначала установите продолжительность движения плота, затем учтите задержку старта и стоянку теплохода.
\end{enumerate}

\clearpage
\section*{Ответы к тренировочному блоку}
\renewcommand{\arraystretch}{1.42}
\begin{center}
\begin{tabularx}{\linewidth}{@{}>{\bfseries\color{accent}}c X@{}}
\toprule
№ & Правильный ответ \\
\midrule
1 & По течению $16$ км/ч, против течения $12$ км/ч; соответственно $3$ ч и $4$ ч; всего $7$ ч. \\
2 & $9$ часов; время встречи~--- 17:30. \\
3 & Скорость лодки в неподвижной воде $15$ км/ч; скорость течения $3$ км/ч. \\
4 & $8$ км/ч. \\
5 & $2$ км/ч. \\
6 & Между городами $80$ км; полный путь $160$ км. \\
7 & $18$ км/ч. \\
8 & $14$ км/ч. \\
9 & $3$ км/ч. \\
10 & $18$ км/ч. \\
\bottomrule
\end{tabularx}
\end{center}
\vspace{11mm}
\noindent\important{Самопроверка.} При подстановке ответа проверьте равенство времён и интервалы отправления. Скорость движения против течения должна оставаться положительной; стоянка не увеличивает расстояние, пройденное судном.
\vfill
\begin{center}
{\small\color{soft} Лёвин Артём Александрович}\par
{\small\color{soft} эксклюзивно для Ярослава Верещагина}
\end{center}
\end{document}
